% META: {"id":"TNSI-STRUCTURES-DONNEES-RE-C03","chapitre":"TNSI-STRUCTURES-DONNEES","type_objet":"remediation","capacites_codes":["C6"],"status":"approved","origin":"generated","status_history":["generated","audited","accepted","approved"],"release_acceptance":"RELEASE_OWNER_BATCH_ACCEPTANCE","acceptance_closure_digest":"sha256:9b3ccf9a81c5520fbb7e03b7b2d3e7bf2057a02834b6d908bf3be5f49f3b553f","acceptance_receipt_digest":"sha256:4fad86f0282e0fa4e0419cca1ad6379ae7af5a280c04a4a90880f90a1c044242"}

%---------------------------------------------------------------
% FICHE DE REMEDIATION — C03A : LIFO/FIFO
%---------------------------------------------------------------

\begin{exercice}{TNSI-STRUCT-RE-C03-EX1}{1}{10}
On empile successivement les valeurs $10$, $20$, $30$ dans une pile, puis on depile deux fois. Quelle est la valeur au sommet de la pile après ces opérations ?
\end{exercice}

% BEGIN-VERIFY
% pile = []
% for v in (10, 20, 30):
%     pile.append(v)
% pile.pop()
% pile.pop()
% assert pile[-1] == 10
% END-VERIFY

\begin{corrige}{TNSI-STRUCT-RE-C03-EX1}
Après avoir empile $10, 20, 30$, la pile est (du bas vers le sommet) $[10, 20, 30]$. Depiler une fois retire $30$ (LIFO), depiler une deuxième fois retire $20$. Il reste $10$ au sommet.
\end{corrige}
